\section*{الفصل \ref{ch:g8:balanced-equations} --- حفظ الكتلة والمعادلات الموزونة}

\begin{solution}{exo:g8:balanced-equations:1}
في \omterm{def:g7:chemical-reactions:reaction}{التفاعل الكيميائي} تساوي الكتلة الكلية للنواتج المتكوّنة الكتلة الكلية
للمتفاعلات المستهلكة.
\end{solution}

\begin{solution}{exo:g8:balanced-equations:2}
المعامل هو 3. ويمثّل 3 \omterm{def:g7:atoms-and-molecules:atom}{ذرات} كربون و$3 \times 2 = 6$
\omterm{def:g7:atoms-and-molecules:atom}{ذرات} أكسجين.
\end{solution}

\begin{solution}{exo:g8:balanced-equations:3}
نعم: 1 C و2 O في كل طرف.
\end{solution}

\begin{solution}{exo:g8:balanced-equations:4}
\ce{2Mg + O2 -> 2MgO}: 2 Mg و2 O في كل طرف.
\end{solution}

\begin{solution}{exo:g8:balanced-equations:5}
النيتروجين: 2 N في اليسار، إذن 2 \ce{NH3}؛ فيصير في اليمين 6 H، إذن 3
\ce{H2}: \ce{N2 + 3H2 -> 2NH3}.
\end{solution}

\begin{solution}{exo:g8:balanced-equations:6}
\omterm{def:g8:balanced-equations:conservation}{بحفظ الكتلة}: $44 - 12 = \qty{32}{g}$ من الأكسجين.
\end{solution}

\begin{solution}{exo:g8:balanced-equations:7}
الغاز المتكوّن يفلت إلى الغرفة، فلا يعود الميزان يزنه. ولم تفنَ أي كتلة:
فالكتلة الناقصة موجودة في \omterm{def:g4:air-a-mixture-of-gases:air}{هواء} الغرفة.
\end{solution}

\begin{solution}{exo:g8:balanced-equations:8}
الكربون: 1 في كل طرف. الهيدروجين: 4 في اليسار، إذن 2 \ce{H2O}. الأكسجين:
$2 + 2 = 4$ في اليمين، إذن 2 \ce{O2}:
\ce{CH4 + 2O2 -> CO2 + 2H2O}.
\end{solution}

\begin{solution}{exo:g8:balanced-equations:9}
تحويل \ce{H2O} إلى \ce{H2O2} يغيّر المادة: فالمادة \ce{H2O2} هي فوق
أكسيد الهيدروجين، لا الماء. ولا يجوز تغيير شيء إلا المعاملات:
\ce{2H2 + O2 -> 2H2O}.
\end{solution}

\begin{solution}{exo:g8:balanced-equations:10}
$20 \times 5 = 100$ \omterm{def:g7:atoms-and-molecules:molecule}{جزيء} أكسجين؛ و$20 \times 4 = 80$ \omterm{def:g7:atoms-and-molecules:molecule}{جزيء}
ماء.
\end{solution}

\begin{solution}{exo:g8:balanced-equations:11}
مع 2 \ce{C2H6}: 4 C، إذن 4 \ce{CO2}؛ و12 H، إذن 6 \ce{H2O}؛ والأكسجين في
اليمين $8 + 6 = 14$، إذن 7 \ce{O2}:
\ce{2C2H6 + 7O2 -> 4CO2 + 6H2O}.
\end{solution}

\begin{solution}{exo:g8:balanced-equations:12}
تبقى القراءة \qty{850.0}{g}. فالصندوق محكم الإغلاق: و\qty{0.4}{g} من الشمع
التي احترقت صارت ثنائي أكسيد الكربون وبخار ماء، وهما لا يزالان داخل
الصندوق مع الأكسجين الذي أخذاه. فلم يدخل شيء ولم يخرج شيء.
\end{solution}

\begin{solution}{pb:g8:balanced-equations:1}
\textbf{1.} من أكسجين \omterm{def:g4:air-a-mixture-of-gases:air}{الهواء}، الذي يتحد بالحديد: فأكسيد الحديد يزن الحديد
مضافًا إليه ذلك الأكسجين.

\textbf{2.} كل شيء يبقى داخل المخبار: فالأكسجين المستهلك كان أصلًا في
المخبار، والأكسيد المتكوّن يبقى فيه. فلا تتغير الكتلة
الكلية.

\textbf{3.} \omterm{def:g8:balanced-equations:conservation}{حفظ الكتلة}: في التفاعل تساوي الكتلة الكلية للنواتج الكتلة
الكلية للمتفاعلات.

\textbf{4.} \ce{4Fe + 3O2 -> 2Fe2O3}.

\textbf{5.} \ce{3Fe + 2O2 -> Fe3O4}.

\textbf{6.} اليسار: 4 Fe، و$3 \times 2 = 6$ O. اليمين: $2 \times 2 = 4$ Fe،
و$2 \times 3 = 6$ O. موزونة.

\textbf{7.} 4 \omterm{def:g7:atoms-and-molecules:atom}{ذرات} حديد لكل 3 \omterm{def:g7:atoms-and-molecules:molecule}{جزيئات} أكسجين، إذن $400 \times
\frac{3}{4} = 300$ \omterm{def:g7:atoms-and-molecules:molecule}{جزيء} أكسجين.

\textbf{8.} $0.28 \times \frac{3}{7} = \qty{0.12}{g}$ من الأكسجين.

\textbf{9.} $0.28 + 0.12 = \qty{0.40}{g}$ من أكسيد الحديد.

\textbf{10.} نعم: يلزم \qty{0.12}{g}، وكان في المخبار نحو
\qty{0.5}{g}.

\textbf{11.} الأكسجين المستهلك غادر \omterm{def:g4:air-a-mixture-of-gases:air}{هواء} المخبار (فهو الآن داخل الأكسيد
الصلب)، فصار في المخبار غاز أقل من قبل، فيدخل إليه \omterm{def:g4:air-a-mixture-of-gases:air}{هواء} الغرفة. وعندئذ
ترتفع قراءة الميزان بمقدار كتلة \omterm{def:g4:air-a-mixture-of-gases:air}{الهواء} الداخل.

\textbf{12.} أُخذ \qty{0.12}{g} من الأكسجين من \omterm{def:g4:air-a-mixture-of-gases:air}{هواء} المخبار.
\end{solution}
